% !TeX program = xelatex
\documentclass[11pt,a4paper]{ctexart}

\usepackage[left=2.4cm,right=2.4cm,top=2.1cm,bottom=2.1cm]{geometry}
\usepackage{amsmath,amssymb,mathtools,braket}
\usepackage{xcolor}
\usepackage[inline]{enumitem}
\usepackage{fancyhdr}
\usepackage[hidelinks]{hyperref}

% ============================================================
% 基本信息
% ============================================================
\newcommand{\coursename}{高等量子力学}
\newcommand{\assignmentnumber}{第 3 次作业}
\newcommand{\teachername}{\underline{\makebox[2.8cm][c]{金彪}}}
\newcommand{\studentname}{\underline{\makebox[2.8cm][c]{林照翔}}}
\newcommand{\studentid}{\underline{\makebox[3.0cm][c]{202628001207046}}}

% ============================================================
% 页面样式
% ============================================================
\setlength{\parindent}{2em}
\setlength{\parskip}{0.25em}
\setlength{\headheight}{14pt}
\pagestyle{fancy}
\fancyhf{}
\lhead{\coursename}
\rhead{\assignmentnumber}
\cfoot{\thepage}
\renewcommand{\headrulewidth}{0.4pt}

% ============================================================
% 列表参数
% ============================================================
\setlist[enumerate,1]{
  label=(\alph*),
  leftmargin=3.2em,
  itemsep=0.55em,
  topsep=0.45em
}

\newlist{inlineenum}{enumerate*}{1}
\setlist[inlineenum,1]{
  label=(\alph*),
  labelsep=0.5em,
  itemjoin={；\quad},
  itemjoin*={；\quad}
}

% ============================================================
% 题目与回答空间
% ============================================================
\newcounter{problem}
\newenvironment{problem}
  {\refstepcounter{problem}%
   \par\medskip\noindent
   \begingroup
   \setlength{\fboxsep}{3pt}%
   \setlength{\fboxrule}{0.8pt}%
   \fbox{\bfseries 题目\chinese{problem}}%
   \endgroup
   \quad\ignorespaces}
  {\par\medskip}

\newenvironment{answer}
  {\par\smallskip\noindent
   \begingroup
   \setlength{\fboxsep}{4pt}%
   \colorbox{black}{\textcolor{white}{\bfseries 回答}}%
   \endgroup
   \quad\ignorespaces}
  {\par\smallskip\noindent\rule{\linewidth}{0.4pt}\medskip}

\newcommand{\answerspace}[1]{\par\vspace*{#1}}

\begin{document}

\begin{center}
  {\LARGE\bfseries \coursename\quad \assignmentnumber}\par
  \vspace{0.7em}
  \small
  教师：\teachername\qquad
  姓名：\studentname\qquad
  学号：\studentid
\end{center}
\vspace{0.4em}

% 本份作业承接原题前十五题，题号从十六开始。
\setcounter{problem}{15}

\begin{problem}
考虑外场作用下的一维谐振子系统：
\begin{equation}
  \hat H=\hat H_0+\hat V,
\end{equation}
其中
\begin{equation}
  \hat H_0=\frac{\hat p^2}{2m}+\frac{m\omega^2\hat x^2}{2},
  \qquad
  \hat V=-\lambda\hat x.
\end{equation}
设初始体系处于相干态上，求 $t$ 时刻的平均值
\begin{equation}
  \overline{x}_t
  ={}_I\Braket{\alpha,t|\hat x_I(t)|\alpha,t}_I
\end{equation}
（准确至 $\lambda^2$ 量级）。
\end{problem}

\begin{answer}
  对于自由谐振子，其哈密顿量 $\hat{H}_0 $ 可写为
  \begin{equation}
    \hat{H}_0
    =\frac{\hat p^2}{2m}+\frac{m\omega^2\hat x^2}{2}
    =\hbar\omega\left(\hat{a}^\dag \hat{a} + 1/2 \right),
  \end{equation}
  其中
  \begin{equation}
    \hat a=\sqrt{\frac{m\omega}{2\hbar}}\left(\hat x+\frac{\mathrm{i}}{m\omega}\hat p\right),\qquad
    \hat a^\dagger=\sqrt{\frac{m\omega}{2\hbar}}\left(\hat x-\frac{\mathrm{i}}{m\omega}\hat p\right).
  \end{equation}
  可以反解出
  \begin{equation}
    \hat{x}
    =\sqrt{\frac{\hbar }{2m\omega } }\left(\hat{a}+\hat{a}^\dag \right).
  \end{equation}
  定义
  \begin{equation}
    \hat{U}_0(t,t_0)
    =\mathrm{e}^{-\mathrm{i}\hat{H}_0(t-t_0)/\hbar},
  \end{equation}
  则相互作用绘景中的算符 $\hat{A}_I(t) $ 定义为
  \begin{equation}
    \hat{A}_I(t)
    =\hat{U}_0^\dag(t,t_0) \hat{A}_S(t) \hat{U}_0(t,t_0).
  \end{equation}
  相互作用绘景下，算符 $\hat{A}_I(t) $ 的时间演化满足如下的微分方程：
  \begin{equation}
    \frac{\mathrm{d}\hat{A}_I(t) }{\mathrm{d}t }
    =\frac{1 }{\mathrm{i}\hbar }\left[\hat{A}_I(t),\hat{H}_0 \right] + \left(\frac{\partial \hat{A}_S(t) }{\partial t}  \right)_I.  
  \end{equation}
  由 $\left[\hat{H}_0,\hat{U}_0(t,t_0) \right]=0 $，分别计算：
  \begin{equation}
    \begin{aligned}
      \left[\hat{H}_0,\hat{a}_I(t) \right]
      &=\left[\hat{H}_0,\hat{U}_0^\dag(t,t_0) \hat{a} \hat{U}_0(t,t_0) \right] \\
      &=\hat{H}_0 \hat{U}_0^\dag(t,t_0) \hat{a} \hat{U}_0(t,t_0) - \hat{U}_0^\dag(t,t_0) \hat{a} \hat{U}_0(t,t_0) \hat{H}_0 \\
      &=\hat{U}_0^\dag(t,t_0) \hat{H}_0 \hat{a} \hat{U}_0(t,t_0) - \hat{U}_0^\dag(t,t_0) \hat{a} \hat{H}_0 \hat{U}_0(t,t_0) \\
      &=\hat{U}_0^\dag(t,t_0) \left[\hat{H}_0,\hat{a} \right] \hat{U}_0(t,t_0)
      =\hat{U}_0^\dag(t,t_0) \left[\hbar\omega\left(\hat{a}^\dag\hat{a}+1/2 \right),\hat{a} \right] \hat{U}_0(t,t_0) \\
      &=-\hbar\omega \hat{U}_0^\dag(t,t_0) \hat{a} \hat{U}_0(t,t_0) \\
      &=-\hbar\omega \hat{a}_I(t),
    \end{aligned}
  \end{equation}
  \begin{equation}
    \begin{aligned}
      \left[\hat{H}_0,\hat{a}_I^\dag(t) \right]
      &=\hat{U}_0^\dag(t,t_0) \left[\hat{H}_0,\hat{a}^\dag \right] \hat{U}_0(t,t_0)
      =\hat{U}_0^\dag(t,t_0) \left[\hbar\omega\left(\hat{a}^\dag\hat{a}+1/2 \right),\hat{a}^\dag \right] \hat{U}_0(t,t_0) \\
      &=\hbar\omega \hat{U}_0^\dag(t,t_0) \hat{a}^\dag \hat{U}_0(t,t_0) \\
      &=\hbar\omega \hat{a}_I^\dag(t).
    \end{aligned}
  \end{equation}
  因此，$\hat{a}_I(t),\hat{a}_I^\dag(t) $ 分别满足：
  \begin{equation}
    \frac{\mathrm{d}\hat{a}_I(t) }{\mathrm{d}t }
    =\frac{1 }{\mathrm{i}\hbar } \left[\hat{a}_I(t),\hat{H}_0 \right] 
    =-\mathrm{i}\omega \hat{a}_I(t),
  \end{equation}
  \begin{equation}
    \frac{\mathrm{d}\hat{a}_I^\dag(t) }{\mathrm{d}t }
    =\frac{1 }{\mathrm{i}\hbar } \left[\hat{a}_I^\dag(t),\hat{H}_0 \right] 
    =\mathrm{i}\omega\hat{a}^\dag_I(t).
  \end{equation}
  设初始时刻 $t_0=0 $，则结合初始条件 $\hat{a}_I(t=0)=\hat{a},\hat{a}_I^\dag(t=0)=\hat{a}^\dag $ 可得：
  \begin{equation}
    \hat{a}_I(t)
    =\mathrm{e}^{-\mathrm{i}\omega t}\hat{a},\quad
    \hat{a}_I^\dag(t)
    =\mathrm{e}^{\mathrm{i}\omega t}\hat{a}^\dag.
  \end{equation}
  于是相互作用绘景中的坐标算符为
  \begin{equation}
    \hat{x}_I(t)
    =\sqrt{\frac{\hbar }{2m\omega } }\left(\hat{a}_I(t)+\hat{a}^\dag_I(t) \right)
    =\sqrt{\frac{\hbar }{2m\omega } } \left(\mathrm{e}^{-\mathrm{i}\omega t}\hat{a} + \mathrm{e}^{\mathrm{i}\omega t}\hat{a}^\dag \right).
  \end{equation}
  相互作用绘景中的微扰算符为
  \begin{equation}
    \hat{V}_I(t)
    =-\lambda\hat{x}_I(t)
    =-\lambda \sqrt{\frac{\hbar }{2m\omega } } \left(\mathrm{e}^{-\mathrm{i}\omega t}\hat{a} + \mathrm{e}^{\mathrm{i}\omega t}\hat{a}^\dag \right).
  \end{equation}
  相互作用绘景中的态矢的定义为
  \begin{equation}
    \Ket{\psi_I(t) }
    =\hat{U}_0^\dag(t,t_0) \Ket{\psi_S(t) },
  \end{equation}
  取 $t=t_0 $ 可得两个绘景中初态的关系：
  \begin{equation}
    \Ket{\psi_I(t_0) }
    =\Ket{\psi_S(t_0) }.
  \end{equation}
  相互作用绘景中态的演化形式上可写为
  \begin{equation}
    \Ket{\psi_I(t) }
    =\hat{U}_I(t,t_0) \Ket{\psi_I(t_0) },\qquad
    \hat{U}_I(t,t_0)
    =\mathcal{T}\exp\left(-\frac{\mathrm{i} }{\hbar }\int_{t_0}^{t}\mathrm{d}t' \hat{V}_I(t') \right),
  \end{equation}
  其中 $\mathcal{T} $ 是编时算符，它把时间参数大的算符放左边。形式上 $\Ket{\alpha,t }_I $ 可写为
  \begin{equation}
    \Ket{\alpha,t }
    =\hat{U}_I(t,0) \Ket{\alpha },
  \end{equation}
  $t $ 时刻坐标算符的期望值为：
  \begin{equation}
    \bar{x}_t
    ={}_I\Braket{\alpha,t|\hat{x}_I(t)|\alpha,t }_I
    =\Braket{\alpha|\hat{U}_I^\dag(t,0)\hat{x}_I(t)\hat{U}_I(t,0)|\alpha }.
  \end{equation}
  把相互作用绘景时间演化算符展开到二阶：
  \begin{equation}
    \begin{aligned}
      \hat{U}_I(t,0)
      &=\mathcal{T}\exp\left(-\frac{\mathrm{i} }{\hbar }\int_{t_0}^{t}\mathrm{d}t' \hat{V}_I(t') \right) \\
      &\approx \hat{I} + \left(-\frac{\mathrm{i} }{\hbar } \int_{0}^{t}\mathrm{d}t_1 \hat{V}_I(t_1)  \right) + \frac{1 }{2! } \left(-\frac{\mathrm{i} }{\hbar }  \right)^2 \int_{0}^{t}\mathrm{d}t_1\int_{0}^{t}\mathrm{d}t_2 \mathcal{T} \left[\hat{V}_I(t_1)\hat{V}_I(t_2) \right] \\
      &\approx \hat{I} + \left(-\frac{\mathrm{i} }{\hbar } \int_{0}^{t}\mathrm{d}t_1 \hat{V}_I(t_1)  \right) +  \left(-\frac{\mathrm{i} }{\hbar }  \right)^2 \int_{0}^{t}\mathrm{d}t_1\int_{0}^{t_1}\mathrm{d}t_2 \hat{V}_I(t_1)\hat{V}_I(t_2) \\
      &\equiv \hat{U}_I^{(0)}(t,0) + \hat{U}_I^{(1)}(t,0) + \hat{U}_I^{(2)}(t,0),
    \end{aligned}
  \end{equation}
  注意到
  \begin{equation}
    \begin{aligned}
      \bar{x}_t
      &=\Braket{\alpha|\hat{U}_I^\dag(t,0)\hat{x}_I(t)\hat{U}_I(t,0)|\alpha } \\
      &\approx \Braket{\alpha|\left(\hat{U}_I^{(0)\dag}(t,0)+\hat{U}_I^{(1)\dag}(t,0)+\hat{U}_I^{(2)\dag}(t,0) \right)\hat{x}_I(t) \left(\hat{U}_I^{(0)}(t,0)+\hat{U}_I^{(1)}(t,0)+\hat{U}_I^{(2)}(t,0) \right)|\alpha } \\
      &\approx \Braket{\hat{x}_I(t) }_\alpha + \Braket{\hat{U}_I^{(1)\dag}(t,0) \hat{x}_I(t) + \hat{x}_I(t)\hat{U}_I^{(1)}(t,0)  }_\alpha \\
      &+\Braket{\hat{U}_I^{(1)\dag}(t,0)\hat{x}_I(t)\hat{U}_I^{(1)}(t,0) + \hat{U}_I^{(2)\dag}(t,0) \hat{x}_I(t) + \hat{x}_I(t)\hat{U}_I^{(2)}(t,0)  }_\alpha \\
      &\equiv \bar{x}_t^{(0)} + \bar{x}_t^{(1)} + \bar{x}_t^{(2)}.
    \end{aligned}
  \end{equation}
  其中
  \begin{equation}
    \begin{aligned}
      \bar{x}_t^{(0)}
      &=\Braket{\hat{x}_I }_\alpha 
      =\sqrt{\frac{\hbar }{2m\omega } } \left(\mathrm{e}^{-\mathrm{i}\omega t}\Braket{\alpha|\hat{a}|\alpha } + \mathrm{e}^{\mathrm{i}\omega t}\Braket{\alpha|\hat{a}^\dag|\alpha } \right) \\
      &=\sqrt{\frac{\hbar }{2m\omega } } \left(\mathrm{e}^{-\mathrm{i}\omega t}\alpha + \mathrm{e}^{\mathrm{i}\omega t}\alpha^* \right).
    \end{aligned}
  \end{equation}
  利用
  \begin{equation}
    \begin{aligned}
      \left[\hat{x}_I(t_1),\hat{x}_I(t) \right]
      =\frac{\hbar }{2m\omega } \left[\mathrm{e}^{-\mathrm{i}\omega t_1}\hat{a} + \mathrm{e}^{\mathrm{i}\omega t_1}\hat{a}^\dag , \mathrm{e}^{-\mathrm{i}\omega t}\hat{a} + \mathrm{e}^{\mathrm{i}\omega t}\hat{a}^\dag \right]
      =\frac{\mathrm{i} \hbar}{m\omega } \sin\left[\omega\left(t-t_1 \right) \right],
    \end{aligned}
  \end{equation}
  有：
  \begin{equation}
    \begin{aligned}
      \bar{x}_t^{(1)}
      &=\Braket{\hat{U}_I^{(1)\dag}(t,0) \hat{x}_I(t) + \hat{x}_I(t)\hat{U}_I^{(1)}(t,0)  }_\alpha
      =\frac{\mathrm{i} }{\hbar } \int_{0}^{t} \mathrm{d}t_1 \Braket{\left[\hat{V}_I(t_1),\hat{x}_I(t) \right] }_\alpha \\
      &=\left(-\lambda \right)\frac{\mathrm{i} }{\hbar } \int_{0}^{t} \mathrm{d}t_1 \Braket{\left[\hat{x}_I(t_1),\hat{x}_I(t) \right] }_\alpha
      =\frac{\lambda }{m\omega } \int_{0}^{t}\mathrm{d}t_1 \sin\left[\omega\left(t-t_1 \right) \right] \\
      &=\frac{\lambda\left(1-\cos \omega t \right) }{m\omega^2 }.
    \end{aligned}
  \end{equation}
  利用
  \begin{equation}
    \begin{aligned}
      \hat{U}_I^{(1)\dag}(t,0) \hat{x}_I(t) \hat{U}_I^{(1)}(t,0)
      &=\frac{1 }{\hbar^2 } \int_{0}^{t}\mathrm{d}t_1\int_{0}^{t}\mathrm{d}t_2 \hat{V}_I(t_1)\hat{x}_I(t)\hat{V}_I(t_2) \\
      &=\frac{1 }{\hbar^2 } \int_{0}^{t}\mathrm{d}t_1\int_{0}^{t_1}\mathrm{d}t_2 \left[\hat{V}_I(t_1)\hat{x}_I(t)\hat{V}_I(t_2) + \hat{V}_I(t_2)\hat{x}_I(t)\hat{V}_I(t_1) \right] \\
    \end{aligned}
  \end{equation}
  有：
  \begin{equation}
    \begin{aligned}
      \bar{x}_t^{(2)}
      &=\Braket{\hat{U}_I^{(1)\dag}(t,0)\hat{x}_I(t)\hat{U}_I^{(1)}(t,0) + \hat{U}_I^{(2)\dag}(t,0) \hat{x}_I(t) + \hat{x}_I(t)\hat{U}_I^{(2)}(t,0)  }_\alpha \\
      &=\frac{1 }{\hbar^2 } \int_{0}^{t}\mathrm{d}t_1\int_{0}^{t_1}\mathrm{d}t_2 \Big\langle\big[\hat{V}_I(t_1)\hat{x}_I(t)\hat{V}_I(t_2) + \hat{V}_I(t_2)\hat{x}_I(t)\hat{V}_I(t_1) \\
      &-\hat{V}_I(t_2)\hat{V}_I(t_1)\hat{x}_I(t) - \hat{x}_I(t)\hat{V}_I(t_1)\hat{V}_I(t_2) \big]\Big\rangle_\alpha \\
      &=-\frac{1 }{\hbar^2 } \int_{0}^{t}\mathrm{d}t_1\int_{0}^{t_1}\mathrm{d}t_2 \Braket{\left[\hat{V}_I(t_2),\left[\hat{V}_I(t_1),\hat{x}_I(t) \right] \right] }_\alpha \\
      &=-\frac{\lambda^2 }{\hbar^2 } \int_{0}^{t}\mathrm{d}t_1\int_{0}^{t_1}\mathrm{d}t_2 \Braket{\left[\hat{x}_I(t_2),\left[\hat{x}_I(t_1),\hat{x}_I(t) \right] \right] }_\alpha \\
      &=0.
    \end{aligned}
  \end{equation}
  综上，准确到二阶小量，
  \begin{equation}
    \bar{x}_t
    \approx \bar{x}_t^{(0)} + \bar{x}_t^{(1)} + \bar{x}_t^{(2)}
    =\sqrt{\frac{\hbar }{2m\omega } } \left(\mathrm{e}^{-\mathrm{i}\omega t}\alpha + \mathrm{e}^{\mathrm{i}\omega t}\alpha^* \right) + \frac{\lambda\left(1-\cos \omega t \right) }{m\omega^2 }.
  \end{equation}
\end{answer}

\begin{problem}
从本征值方程
\begin{equation}
  \hat x_H(t)\Ket{x,t}_H=x\Ket{x,t}_H
\end{equation}
出发，导出自由粒子和谐振子系统中
\begin{equation}
  K^*(x,t;x_0,t_0)={}_H\Braket{x_0,t_0|x,t}_H
\end{equation}
满足的微分方程，进而导出传播子 $K(x,t;x_0,t_0)$ 的表达式。
\end{problem}

\begin{answer}
  由海森堡绘景中算符、态矢的定义，有：
  \begin{equation}
    \hat{x}_H(t) = \hat{U}^\dag(t,t_0) \hat{x} \hat{U}(t,t_0),\qquad
    \Ket{x,t }_H = \hat{U}^\dag(t,t_0) \Ket{x },
  \end{equation}
  当 $t=t_0 $ 时，有：
  \begin{equation}
    \hat{x}_H(t_0) = \hat{x},\qquad
    \Ket{x,t_0 }_H = \Ket{x }.
  \end{equation}
  于是可以在海森堡绘景下表达传播子：
  \begin{equation}
    K(x,t;x_0,t_0)
    =\Braket{x|\hat{U}(t,t_0)|x_0 }
    ={}_H\Braket{x,t|x_0,t_0 }_H,\qquad
    K^*(x,t;x_0,t_0)
    ={}_H\Braket{x_0,t_0|x,t }_H.
  \end{equation}
  对于不显含时间的哈密顿量，时间演化算符形式上可写为
  \begin{equation}
    \hat{U}(t,t_0) = \mathrm{e}^{-\mathrm{i}\hat{H}(t-t_0)/\hbar},\qquad
    \Ket{x,t }_H = \hat{U}^\dag(t,t_0) \Ket{x } = \mathrm{e}^{\mathrm{i}\hat{H}(t-t_0)/\hbar} \Ket{x }.
  \end{equation}
  两边对 $t $ 求导，得：
  \begin{equation}
    -\mathrm{i}\hbar\frac{\partial }{\partial t } \Ket{x,t }_H
    =\hat{H} \Ket{x,t }_H.
  \end{equation}
  上式两边左乘 ${}_H\Bra{x_0,t_0 } $ 得到传播子满足的微分方程：
  \begin{equation}
    -\mathrm{i}\hbar\frac{\partial }{\partial t}K^*(x,t;x_0,t_0)
    ={}_H\Braket{x_0,t_0|\hat{H}|x,t }_H. 
    \label{eq:time-evolution}
  \end{equation}
  \begin{enumerate}
    \item 对于自由粒子，其哈密顿量为 $\hat{H}=\hat{p}^2/(2m) $。海森堡绘景中不显含时间的算符演化满足
    \begin{equation}
      \frac{\mathrm{d} \hat{A}_H(t)}{\mathrm{d}t }
      =\frac{1 }{\mathrm{i}\hbar } \left[\hat{A}_H(t),\hat{H} \right].
    \end{equation}
    \begin{equation}
      \begin{aligned}
        \frac{\mathrm{d} \hat{x}_H(t)}{\mathrm{d}t }
        =\frac{1 }{\mathrm{i}\hbar } \left[\hat{x}_H(t),\hat{H} \right]
        =\frac{1 }{\mathrm{i}\hbar } \hat{U}^\dag(t,t_0) \left[\hat{x},\hat{H} \right] \hat{U}(t,t_0)
        =\hat{U}^\dag(t,t_0)  \frac{\hat{p} }{m } \hat{U}(t,t_0)
        =\frac{\hat{p}_H(t) }{m } ,
      \end{aligned}
    \end{equation}
    \begin{equation}
      \begin{aligned}
        \frac{\mathrm{d} \hat{p}_H(t)}{\mathrm{d}t }
        =\frac{1 }{\mathrm{i}\hbar } \left[\hat{p}_H(t),\hat{H} \right]
        =\frac{1 }{\mathrm{i}\hbar } \hat{U}^\dag(t,t_0) \left[\hat{p},\hat{H} \right] \hat{U}(t,t_0)
        =0,
      \end{aligned}
    \end{equation}
    结合初始条件 $\hat{p}_H(t_0)=\hat{p},\hat{x}_H(t_0)=\hat{x} $ 解得：
    \begin{equation}
      \hat{p}_H(t) = \hat{p}_H(t_0) = \hat{p},\qquad
      \hat{x}_H(t) = \hat{x}_H(t_0) + \frac{t-t_0 }{m } \hat{p}_H(t_0) = \hat{x} + \frac{t-t_0 }{m } \hat{p}.  
    \end{equation}
    把上式代入本征方程 $\hat{x}_H(t) \Ket{x,t }_H = x \Ket{x,t }_H $，并左乘 ${}_H\Bra{x_0,t_0 } $ 得到：
    \begin{equation}
      \left(x_0 - \frac{\mathrm{i}\hbar\left(t-t_0 \right) }{m } \frac{\partial }{\partial x_0 }  \right) K^*(x,t;x_0,t_0)
      =x K^*(x,t;x_0,t_0).
    \end{equation}
    可以解得：
    \begin{equation}
      K^*(x,t;x_0,t_0)
      =C_1(x,t;t_0) \exp\left[-\frac{\mathrm{i}m\left(x-x_0 \right)^2 }{2\hbar (t-t_0) }  \right]
    \end{equation}
    再计算矩阵元
    \begin{equation}
      \begin{aligned}
        {}_H\Braket{x_0,t_0|\hat{H}|x,t }_H
        &=\Braket{x_0|\frac{\hat{p}^2 }{2m } |x,t }_H
        =-\frac{\hbar^2 }{2m }\frac{\partial^2 }{\partial x_0^2 } \Braket{x_0|x,t }_H \\
        &=-\frac{\hbar^2 }{2m }\frac{\partial^2 }{\partial x_0^2 } {}_H\Braket{x_0,t_0|x,t }_H 
        =-\frac{\hbar^2 }{2m }\frac{\partial^2 }{\partial x_0^2 } K^*(x,t;x_0,t_0),
      \end{aligned}
    \end{equation}
    代入 \eqref{eq:time-evolution} 得到方程：
    \begin{equation}
      -\mathrm{i}\hbar\frac{\partial }{\partial t}K^*(x,t;x_0,t_0)
      =-\frac{\hbar^2 }{2m }\frac{\partial^2 }{\partial x_0^2 } K^*(x,t;x_0,t_0),
    \end{equation}
    也即：
    \begin{equation}
      -\mathrm{i} \hbar \frac{\partial C_1(x,t;t_0) }{\partial t }
      =\frac{\mathrm{i}\hbar }{2(t-t_0) } C_1(x,t;t_0),  
    \end{equation}
    解得：
    \begin{equation}
      C_1(x,t;t_0)
      =C_2(x,t_0) \frac{1 }{\sqrt{t-t_0} } ,
    \end{equation}
    \begin{equation}
        K^*(x,t;x_0,t_0)
        =\frac{C_2(x,t_0) }{\sqrt{t-t_0} } \exp\left[-\frac{\mathrm{i}m\left(x-x_0 \right)^2 }{2\hbar (t-t_0) }  \right].
    \end{equation}
    利用初始条件 $\displaystyle{\lim_{t\to t_0^+} K(x,t;x_0,t_0) = \delta(x-x_0) } $ 以及公式
    \begin{equation}
      \delta(x-x_0)
      =\lim_{\tau\to 0^+} \sqrt{\frac{m }{-2\pi\mathrm{i}\hbar\tau } }\exp\left[-\frac{\mathrm{i}m(x-x_0)^2 }{2\hbar\tau }  \right],
      \label{eq:delta}
    \end{equation}
    对比可得：
    \begin{equation}
      C_2(x,t_0)
      =\sqrt{\frac{m }{-2\pi\mathrm{i}\hbar } },\qquad
      K^*(x,t;x_0,t_0)
      =\sqrt{\frac{m }{-2\pi\mathrm{i}\hbar(t-t_0) } } \exp\left[-\frac{\mathrm{i}m\left(x-x_0 \right)^2 }{2\hbar (t-t_0) }  \right],
    \end{equation}
    \begin{equation}
      K(x,t;x_0,t_0)
      =\sqrt{\frac{m }{2\pi\mathrm{i}\hbar(t-t_0) } } \exp\left[\frac{\mathrm{i}m\left(x-x_0 \right)^2 }{2\hbar (t-t_0) }  \right].
    \end{equation}
    \item 对于谐振子，哈密顿量为 $\hat{H}=\hat{p}^2/(2m) + m\omega^2 \hat{x}^2/2 $，不显含时间，因此
    \begin{equation}
      \begin{aligned}
        \frac{\mathrm{d}\hat{x}_H(t) }{\mathrm{d}t }
        =\frac{1 }{\mathrm{i}\hbar }\left[\hat{x}_H(t),\hat{H} \right]  
        =\frac{\hat{p}_H(t) }{m },
      \end{aligned}
    \end{equation}
    \begin{equation}
      \begin{aligned}
        \frac{\mathrm{d} \hat{p}_H(t)}{\mathrm{d}t }
        =\frac{1 }{\mathrm{i}\hbar }\left[\hat{p}_H(t),\hat{H} \right]
        =-m\omega^2 \hat{x}_H(t),
      \end{aligned}
    \end{equation}
    联立得到
    \begin{equation}
      \ddot{\hat{x}}_H(t)  + \omega^2 \hat{x}_H(t)
      =0.
    \end{equation}
    结合初始条件
    \begin{equation}
      \hat{x}_H(t_0)  = \hat{x},\qquad
      \hat{p}_H(t_0) = \hat{p},\qquad
      \dot{\hat{x}}_H(t_0) = \frac{\hat{p}_H(t_0) }{m } = \frac{\hat{p} }{m },
    \end{equation}
    解得
    \begin{equation}
      \hat{x}_H(t) = \hat{x} \cos \left[\omega(t-t_0) \right] + \frac{\hat{p} }{m\omega } \sin\left[\omega(t-t_0) \right],
    \end{equation}
    \begin{equation}
      \hat{p}_H(t) = -m\omega \hat{x} \sin\left[\omega(t-t_0) \right] + \hat{p} \cos\left[\omega(t-t_0) \right].
    \end{equation}
    上式代入海森堡绘景坐标算符本征方程：
    \begin{equation}
      \left(\hat{x} \cos \left[\omega(t-t_0) \right] + \frac{\hat{p} }{m\omega } \sin\left[\omega(t-t_0) \right] \right) \Ket{x,t }_H
      =x\Ket{x,t }_H,
    \end{equation}
    左乘 ${}_H\Bra{x_0,t_0 } $ 得：
    \begin{equation}
      \left(x_0\cos\left[\omega(t-t_0) \right] - \frac{\mathrm{i}\hbar }{m\omega }\sin\left[\omega(t-t_0) \right] \frac{\partial }{\partial x_0 }  \right) K^*(x,t;x_0,t_0)
      =x K^*(x,t;x_0,t_0),
    \end{equation}
    整理得到：
    \begin{equation}
      \frac{\partial K^*(x,t;x_0,t_0) }{\partial x_0 }
      =\frac{\mathrm{i}m\omega }{\hbar } \left(x\csc\left[\omega(t-t_0) \right] - x_0\cot\left[\omega(t-t_0) \right] \right) K^*(x,t;x_0,t_0),
    \end{equation}
    解得：
    \begin{equation}
      K^*(x,t;x_0,t_0)
      =C_1(x,t;t_0) \exp\left[\frac{\mathrm{i}m\omega }{\hbar } \left(-\frac{x_0^2 \cot\left[\omega(t-t_0) \right] }{2 } + x x_0 \csc\left[\omega(t-t_0) \right] \right) \right],
    \end{equation}
    再计算矩阵元
    \begin{equation}
      \begin{aligned}
        {}_H\Braket{x_0,t_0|\hat{H}|x,t }_H
        &=\Braket{x_0|\frac{\hat{p}^2 }{2m } + \frac{1 }{2 } m\omega^2\hat{x}^2 |x,t }_H
        =\left(-\frac{\hbar^2 }{2m }\frac{\partial^2 }{\partial x_0^2 } + \frac{m\omega^2 x_0^2 }{2 } \right) \Braket{x_0|x,t }_H \\
        &=\left(-\frac{\hbar^2 }{2m }\frac{\partial^2 }{\partial x_0^2 } + \frac{m\omega^2 x_0^2 }{2 } \right) {}_H\Braket{x_0,t_0|x,t }_{H} \\
        &=\left(-\frac{\hbar^2 }{2m }\frac{\partial^2 }{\partial x_0^2 } + \frac{m\omega^2 x_0^2 }{2 } \right) K^*(x,t;x_0,t_0),
      \end{aligned}
    \end{equation}
    代入 \eqref{eq:time-evolution} 得到谐振子传播子满足的微分方程
    \begin{equation}
      -\mathrm{i}\hbar\frac{\partial }{\partial t}K^*(x,t;x_0,t_0)
      =\left(-\frac{\hbar^2 }{2m }\frac{\partial^2 }{\partial x_0^2 } + \frac{m\omega^2 x_0^2 }{2 } \right) K^*(x,t;x_0,t_0).
    \end{equation}
    整理得到：
    \begin{equation}
      -\mathrm{i}\hbar \frac{\partial C_1(x,t;t_0) }{\partial t }
      =\left(\frac{\mathrm{i}\hbar\omega }{2 } \cot\left[\omega(t-t_0) \right] + \frac{m\omega^2 x^2 }{2 }\csc^2\left[\omega(t-t_0) \right]  \right) C_1(x,t;t_0), 
    \end{equation}
    利用积分公式
    \begin{equation}
      \int \cot(\omega\tau) \mathrm{d}\tau 
      = \frac{1 }{\omega } \ln\left|\sin(\omega\tau) \right|+C,\quad
      \int \csc^2(\omega\tau)\mathrm{d}\tau
      =-\frac{1 }{\omega } \cot(\omega\tau) + C,
    \end{equation}
    解得：
    \begin{equation}
      C_1(x,t;t_0)
      =\frac{C_2(x,t_0) }{\sqrt{\sin\left[\omega(t-t_0) \right]} } \exp\left[-\frac{\mathrm{i}m\omega x^2 }{2\hbar } \cot\left[\omega(t-t_0) \right] \right],
    \end{equation}
    \begin{equation}
      K^*(x,t;x_0,t_0)
      =\frac{C_2(x,t_0) }{\sqrt{\sin\left[\omega(t-t_0) \right]} } \exp\left[-\frac{\mathrm{i}m\omega \left(\left(x^2+x_0^2 \right)\cos\left[\omega(t-t_0) \right]-2x x_0 \right)}{2\hbar\sin\left[\omega(t-t_0) \right] }  \right],
    \end{equation}
    考虑 $t\to t_0^+ $，此时 $\sin\left[\omega(t-t_0) \right]\sim \omega(t-t_0),\cos\left[\omega(t-t_0) \right]\sim 1 $，于是
    \begin{equation}
      \lim_{t\to t_0^+} K^*(x,t;x_0,t_0)
      =\lim_{t\to t_0^+}\frac{C_2(x,t_0) }{\sqrt{\omega(t-t_0) } } \exp\left[-\frac{\mathrm{i}m \left(x-x_0 \right)^2}{2\hbar(t-t_0) }  \right]
    \end{equation}
    由初始条件 $\displaystyle{\lim_{t\to t_0^+}K(x,t;x_0,t_0)=\delta(x-x_0) }$ 以及 \eqref{eq:delta} 可得：
    \begin{equation}
      C_2(x,t_0)
      =\sqrt{\frac{m\omega }{-2\pi\mathrm{i}\hbar } },
    \end{equation}
    \begin{equation}
      K(x,t;x_0,t_0)
      =\sqrt{\frac{m\omega }{2\pi\mathrm{i}\hbar \sin\left[\omega(t-t_0) \right] } } \exp\left[\frac{\mathrm{i}m\omega \left(\left(x^2+x_0^2 \right)\cos\left[\omega(t-t_0) \right]-2x x_0 \right)}{2\hbar\sin\left[\omega(t-t_0) \right] }  \right].
    \end{equation}
  \end{enumerate}
\end{answer}

\begin{problem}
设自由粒子的初始态为
\begin{equation}
  \psi(x,t_0)
  =\frac{1}{\left(2\pi a^2\right)^{1/4}}
  \exp\left[-\frac{(x-x_0)^2}{4a^2}\right].
\end{equation}
试通过传播子求 $t$ 时刻的波函数 $\psi(x,t)$，并讨论
$\left|\psi(x,t)\right|^2$ 在 $t$ 值很大时的渐进行为。
\end{problem}

\begin{answer}
  自由粒子传播子为
    \begin{equation}
    K(x,t;x',t_0)
    =\sqrt{\frac{m }{2\pi\mathrm{i}\hbar(t-t_0) } } \exp\left[\frac{\mathrm{i}m\left(x-x' \right)^2 }{2\hbar (t-t_0) }  \right],
  \end{equation}
  $t $ 时刻波函数为
  \begin{equation}
    \begin{aligned}
      \psi(x,t)
      &=\int_{-\infty}^{+\infty} K(x,t;x',t_0) \psi(x',t_0)\mathrm{d}x' \\
      &=\frac{1}{\left(2\pi a^2\right)^{1/4}}\sqrt{\frac{m }{2\pi\mathrm{i}\hbar(t-t_0) } }\int_{-\infty}^{+\infty} \exp\left[\frac{\mathrm{i}m\left(x-x' \right)^2 }{2\hbar (t-t_0) } - \frac{(x'-x_0)^2}{4a^2} \right]\mathrm{d}x' \\
      &=\frac{1}{\left(2\pi a^2\right)^{1/4}}\sqrt{\frac{m }{2\pi\mathrm{i}\hbar(t-t_0) } }\int_{-\infty}^{+\infty} \exp\bigg[\left(-\frac{1 }{4a^2 }+\frac{\mathrm{i}m }{2\hbar(t-t_0) }   \right)x'^2 + \left(\frac{x_0 }{2a^2 }-\frac{\mathrm{i}mx }{\hbar(t-t_0) } \right)x' \\
      &+\frac{\mathrm{i}mx^2 }{2\hbar(t-t_0) }-\frac{x_0^2 }{4a^2 }   \bigg]\mathrm{d}x',
    \end{aligned}
  \end{equation}
  定义无量纲时间参数
  \begin{equation}
    \eta
    =\frac{\hbar(t-t_0) }{2ma^2 }, 
  \end{equation}
  \begin{equation}
    \psi(x,t)
    =\frac{1}{\left(2\pi a^2\right)^{1/4}}\sqrt{\frac{1 }{4\pi\mathrm{i}a^2 \eta} } \exp\left(\frac{\mathrm{i}x^2-\eta x_0^2 }{4a^2\eta }  \right) \int_{-\infty}^{+\infty} \exp\left[\left(-\frac{\eta-\mathrm{i} }{4a^2\eta }   \right)x'^2 + \left(\frac{\eta x_0 - \mathrm{i}x }{2a^2\eta } \right)x' \right]\mathrm{d}x'.
  \end{equation}
  利用高斯积分公式
  \begin{equation}
    \int_{-\infty}^{+\infty} \mathrm{e}^{-Ax^2+Bx}\mathrm{d}x
    =\sqrt{\frac{\pi }{A } }\exp\left(\frac{B^2 }{4A }  \right),\quad
    A=\frac{\eta-\mathrm{i} }{4a^2\eta },\quad
    B=\frac{\eta x_0 - \mathrm{i}x }{2a^2\eta },
  \end{equation}
  \begin{equation}
    \begin{aligned}
      \psi(x,t)
      &=\frac{1}{\left(2\pi a^2\right)^{1/4}} \frac{1 }{\sqrt{1+\mathrm{i}\eta} }  \exp\left[-\frac{(x-x_0)^2 }{4a^2(1+\mathrm{i}\eta) }  \right],
    \end{aligned}
  \end{equation}
  \begin{equation}
    \left|\psi(x,t) \right|^2
    =\frac{1 }{\left(2\pi a^2 \right)^{1/2} } \frac{1 }{\sqrt{1+\eta^2} } \exp\left[-\frac{(x-x_0)^2 }{2\left(1+\eta^2 \right)a^2 }  \right].
  \end{equation}
  当 $t\to +\infty $ 时，$\eta\to +\infty $，概率密度仍是高斯波包，波包中心不变，但展宽随时间线性增大。

\end{answer}

\begin{problem}
设系统经典拉氏量为
\begin{equation}
  L=\frac{m\dot{x}^2}{2}-\lambda x.
\end{equation}
求传播子 $K(x,t;x_0,t_0)$ 的表达式。
\end{problem}

\begin{answer}
\answerspace{12cm}
\end{answer}

\begin{problem}
动量表象传播子 $k(p,t;p_0,t_0)$ 定义为
\begin{equation}
  \varphi(p,t)
  =\int k(p,t;p_0,t_0)\varphi(p_0,t_0)\,\mathrm{d}p_0.
\end{equation}
已知动量表象波函数 $\varphi(p,t)$ 与坐标表象波函数 $\psi(x,t)$ 间的变换关系为
\begin{equation}
  \psi(x,t)
  =\frac{1}{\sqrt{2\pi\hbar}}
  \int \mathrm{e}^{\mathrm{i}px/\hbar}\varphi(p,t)\,\mathrm{d}p.
\end{equation}
试求自由粒子系统的动量表象传播子。
\end{problem}

\begin{answer}
\answerspace{11cm}
\end{answer}

\end{document}
